\section{Introduction}
\label{sec:introduction}
A scientific-agent benchmark that marks a conclusion ``correct'' makes three
claims at once. The conclusion agrees with how the data were generated, the
agent's experiment supports it, and the grader credited it for that reason.
These are different measurement objects. A conclusion can be true but
unsupported, as when an agent ignores a significant confirmatory result and
the effect happens to be absent. It can also be supported but untrue, as when
the confirmation sample contains a chance association. Benchmarks for
data-driven discovery
\citep{chen2025scienceagentbench,majumder2025discoverybench,wang2026firebench,lu2024aiscientist}
grade conclusions through pipelines that execute agent code, hold out data,
and check reported results. We ask: \emph{when truth, evidence, and credit
disagree, which disagreements does a pipeline reward, and which pipeline
faults create new ones?} Controlled hypothesis-testing tasks with known
generating processes make all three quantities observable.

Our contributions are:
\begin{itemize}
\item a framework that scores truth, evidence, decision consistency, and
credit separately, with a frozen oracle that recomputes confirmatory evidence
independently of the agent;
\item an instrumented audit of a 1,680-run experiment in which agents wrote
their own confirmation scripts. 42\% of single confirmation calls fit more
than one candidate on held-out data, and 11 of the 19 submissions credited
only by an answer-only grader were not evidence violations;
\item a protocol in which the harness runs one locked test and decisions cite
an immutable evidence record, with single-boundary faults logged for
manifestation, evaluated on 40 new tasks.
\end{itemize}

\section{Related Work}
\label{sec:related}
\citet{kapoor2024agents} and \citet{zhu2025abc} catalog validity flaws in
agent benchmarks. \citet{shao2026protocol} define protocol validity as the
intended capability remaining necessary for success, and audit shortcut
exposure across 15 benchmarks. \citet{chen2026unreliable} treats an agent
harness as software under test, using oracle checks, injected defects, and
replay to correct labels. \citet{he2026sscbench} show that a fault-injection
result supports a claim about an agent only if the faulty observation and the
counter-evidence actually became visible, measured per execution. We
therefore claim no novelty for fault injection, replay, or the observation
that dropping a grading check inflates scores. We add the scientific-evidence
dimension: with a known generating process and an independent oracle,
\emph{correct} and \emph{supported} become separately measurable. Following
\citet{he2026sscbench}, we report whether each fault triggered and reached the
agent. Among scientific-agent benchmarks
\citep{chan2025mlebench,starace2025paperbench,siegel2024corebench},
FIRE-Bench \citep{wang2026firebench} asks for evidence-backed conclusions. We
ask whether the instrument can tell such conclusions from merely correct ones.
Preregistration and holdout discipline
\citep{nosek2018preregistration,dwork2015reusable} guard against analytic
flexibility \citep{simmons2011false,gelman2014statistical}. Our protocol
enforces them in the harness, consistent with treating validity as a property
of the inferences drawn from a score \citep{messick1995validity}.
